\documentclass[../../../include/open-logic-section]{subfiles}

\begin{document}
\olfileid{sth}{ordinals}{zfm}	
\olsection{$\ZFminus$: ఒక మైలురాయి}

ప్రతిస్థాపనను ఎలా
సమర్థించాలి (అసలు
సమర్థించగలమా) అనే
ప్రశ్నకు సరళమైన
జవాబు లేదు. అందువల్ల
దానిని \olref[replacement][]{chap}కు
వదిలేస్తాం. అయితే
ప్రతిస్థాపనను చేర్చడంతో
మరొక ముఖ్యమైన
మైలురాయిని చేరుకున్నాం.
ఇప్పుడు $\ZFminus$
సిద్ధాంతానికి అవసరమైన
స్వీకృతాలన్నీ మన
వద్ద ఉన్నాయి. వివరంగా:

\begin{defn}
$\ZFminus$ సిద్ధాంతంలో
ఈ స్వీకృతాలు ఉంటాయి:
సమితుల సమానత్వ సూత్రం,
సమ్మేళనం, జంటలు,
ఘాత సమితులు, అనంతత్వం,
వేరుచేయడం మరియు
ప్రతిస్థాపన పథకాల
అన్ని సందర్భాలు.
మరోలా చెప్పాలంటే,
$\Zminus$కు ప్రతిస్థాపనను
జోడిస్తే $\ZFminus$.
\end{defn}
\noindent
ఇది \emph{సెర్మెలో--ఫ్రెంకెల్}
సమితి సిద్ధాంతాన్ని
సూచిస్తుంది (తరువాత
చర్చించబోయే ఒకదాన్ని
\emph{తీసివేసిన} రూపం).
\citeyear{Fraenkel1922}లో
ప్రతిస్థాపనను
రూపొందించిన ఘనత
ఫ్రెంకెల్‌కు ఆపాదించబడుతుంది
కాబట్టి అతనికి ఆ
గౌరవం లభించింది;
అయితే తొలి కచ్చితమైన
రూపకల్పన
\citet{Skolem1922}ది.

\end{document}
